\chapter{Conformité et abus de marché}\label{ch:fm:compliance-and-market-abuse}

En août 2022, la Financial Conduct Authority britannique a infligé à Citigroup Global Markets une amende de 12.6~millions de livres parce que, lorsque la Market Abuse Regulation
a étendu en 2016 l'obligation de surveiller les ordres autant que les transactions, la société ne l'a pas correctement mise en œuvre~: il lui a fallu dix-huit mois pour identifier et
évaluer les risques d'\omterm{def:fm:compliance-and-market-abuse:abuse}{abus de marché} auxquels son activité était exposée, et, dans l'intervalle, elle ne pouvait pas surveiller efficacement son trading pour certains types
d'\omterm{def:fm:compliance-and-market-abuse:abuse}{opérations d'initié} et de manipulation de marché. Personne n'a été accusé d'\omterm{def:fm:compliance-and-market-abuse:abuse}{abus de marché}. La défaillance tenait à des lacunes importantes dans la surveillance de la société~: pour certaines formes
d'abus, aucune alerte n'aurait été générée, si bien que personne n'aurait pu en lire une. La conformité est une fonction qui doit voir le trading avant de pouvoir
la juger.

\section{Le règlement de marché auquel est soumis un trader}

La conduite d'un trader est délimitée par deux interdictions et par les obligations qui les appuient. Les interdictions portent sur l'utilisation d'informations que d'autres n'ont pas
et sur la distorsion des prix~; le livre~9 décrit les schémas de manipulation (spoofing, layering, manipulation du cours de clôture, déclenchement de momentum, manipulation de
l'indice de référence) et le wash trading du livre~3. Ce chapitre porte sur la fonction qui permet à une société de respecter ces interdictions.

\begin{definition}[Abus de marché, information privilégiée, opération d'initié]\label{def:fm:compliance-and-market-abuse:abuse}
L'\emph{abus de marché}\index{abus de marché} recouvre les opérations d'initié, la divulgation illicite d'une information privilégiée et la manipulation de marché, y compris
les tentatives. L'\emph{information privilégiée}\index{information privilégiée} est une information de nature précise, non rendue publique, qui concerne directement ou indirectement
à un émetteur ou à un instrument financier, et qui, rendue publique, serait susceptible d'avoir un effet significatif sur les cours. \emph{Opération
d'initié}\index{opération d'initié} consiste à utiliser une information privilégiée pour acheter ou vendre, pour soi-même ou pour autrui, les instruments auxquels elle se rapporte, y compris
en annulant ou en modifiant un ordre passé avant la réception de l'information.
\end{definition}

\begin{dated}{2026-09}{La Market Abuse Regulation de l'UE}\label{dat:fm:compliance-and-market-abuse:mar}
La Regulation (EU) No 596/2014 définit l'\omterm{def:fm:compliance-and-market-abuse:abuse}{information privilégiée} (article~7) et l'\omterm{def:fm:compliance-and-market-abuse:abuse}{opération d'initié} (article~8)~; elle interdit les \omterm{def:fm:compliance-and-market-abuse:abuse}{opérations d'initié}, le fait de recommander ou
d'inciter à en réaliser et la divulgation illicite d'\omterm{def:fm:compliance-and-market-abuse:abuse}{information privilégiée} (article~14), ainsi que la manipulation de marché (article~15)~; elle encadre les sondages de marché, dont les
enregistrements sont conservés pendant au moins cinq ans (article~11)~; elle impose à toute personne qui organise ou exécute professionnellement des transactions de maintenir
des dispositifs, systèmes et procédures efficaces pour détecter et signaler les ordres et transactions suspects, et pour notifier à l'autorité compétente
sans délai en cas de soupçon raisonnable (article~16(2))~; elle exige des émetteurs et de ceux qui agissent pour leur compte qu'ils tiennent des \omterm{def:fm:compliance-and-market-abuse:barrier}{listes d'initiés} (article~18).
\end{dated}

\begin{dated}{2026-09}{Règles américaines sur les opérations d'initié et les transactions personnelles}\label{dat:fm:compliance-and-market-abuse:us}
La Rule 10b-5 (17 CFR 240.10b-5) rend illicite le fait d'employer tout dispositif, tout stratagème ou tout artifice pour frauder, ou de se livrer à tout acte qui s'apparente à une
fraude, dans le cadre de l'achat ou de la vente de toute valeur mobilière~; les poursuites pour \omterm{def:fm:compliance-and-market-abuse:abuse}{opérations d'initié} reposent sur cette règle. La section~15(g) de l'Exchange Act (15
U.S.C. 78o(g)) impose à tout broker-dealer enregistré d'établir, de maintenir et de faire appliquer des politiques et procédures écrites raisonnablement conçues pour
prévenir l'utilisation abusive d'\omterm{def:fm:compliance-and-market-abuse:abuse}{information privilégiée} (material non-public information). La Rule 204A-1 (17 CFR 275.204A-1) exige que le code de déontologie d'un conseiller en investissement enregistré
oblige ses personnes ayant accès à l'information à déclarer leurs avoirs et, dans les 30 jours suivant chaque trimestre, leurs \omterm{def:fm:compliance-and-market-abuse:pad}{transactions personnelles}, et à obtenir
une \omterm{def:fm:compliance-and-market-abuse:pad}{autorisation préalable} avant tout achat lors d'une introduction en bourse ou d'une offre restreinte.
\end{dated}

\section{La surveillance en tant que fonction}

Le livre~15 a construit le système de surveillance des transactions~: des détecteurs qui notent ordres et transactions, des alertes de surveillance, le tri des alertes et l'archive des communications.
La fonction conformité prend en charge la suite~: quels détecteurs tournent sur quelles activités, où se situent leurs seuils, qui examine les alertes, et
ce qui est déclaré.

\begin{definition}[Déclaration d'opérations et d'ordres suspects]\label{def:fm:compliance-and-market-abuse:stor}
Une \emph{déclaration d'opérations et d'ordres suspects}\index{déclaration d'opérations et d'ordres suspects} (STOR) est la notification qu'une société qui organise ou
exécute des transactions doit adresser à son régulateur, sans délai, lorsqu'elle a un soupçon raisonnable qu'un ordre ou une transaction puisse relever d'une \omterm{def:fm:compliance-and-market-abuse:abuse}{opération
d'initié} ou d'une manipulation de marché, tentée ou accomplie.
\end{definition}

L'échec de l'accroche était un problème de couverture~: des détecteurs qui n'ont pas été lancés ne peuvent pas alerter. L'échec plus subtil tient à la capacité. Tout détecteur a un seuil~; en dessous
de ce seuil rien n'est examiné, au-dessus chaque alerte attend un analyste. Si on le place bas, la file déborde, et les analystes qui examinent les alertes dans l'ordre où elles arrivent
ne voient surtout que du bruit~; si on le place haut, les vrais cas ne déclenchent jamais d'alerte.

\begin{proposition}[Remplir la file]\label{prop:fm:compliance-and-market-abuse:queue}
Soit un détecteur qui note $N$ candidats, et soit un taux d'alerte $q$ qui envoie à l'examen les $qN$ candidats les mieux notés, parmi lesquels $\mathrm{TP}(q)$ sont de vrais cas,
non décroissant en $q$, avec une précision $\mathrm{TP}(q)/(qN)$ non croissante en $q$. Si les analystes examinent au plus $C$ alertes, dans l'ordre d'arrivée, le
nombre espéré de vrais cas trouvés vaut $\mathrm{TP}(q)$ pour $q\le C/N$ et $C\,\mathrm{TP}(q)/(qN)$ au-delà, et il est maximal en $q^*=C/N$~: le
seuil dont les alertes remplissent tout juste la capacité des analystes.
\end{proposition}

\begin{proof}
En dessous de $C/N$, toutes les alertes sont examinées et les vrais cas trouvés valent $\mathrm{TP}(q)$, quantité qui ne diminue pas quand $q$ augmente. Au-dessus, une part $C/(qN)$ des
alertes est examinée, dans un ordre indépendant du score, de sorte que le nombre espéré de vrais cas trouvés vaut $C$ fois la précision, qui ne croît pas avec
$q$. Les deux morceaux valent tous deux $\mathrm{TP}(C/N)$ en $q=C/N$.
\end{proof}

\omcode{code/firm/compliance/firm_compliance.py}{76}{93}{La file de tri~: alertes au-dessus du seuil pour un taux d'alerte choisi, examinées dans la limite de
la capacité, dans l'ordre d'arrivée ou dans l'ordre des scores, et le taux d'alerte qui trouve le plus de vrais cas.}\label{lst:fm:compliance-and-market-abuse:triage}

\section{Tutoriel~: la file d'alertes et le desk d'autorisation préalable}\label{tut:fm:compliance-and-market-abuse:tut}

\begin{tutorial}
\textbf{Objectif.} Traiter une année d'\omterm{def:fm:compliance-and-market-abuse:pad}{autorisations préalables} sur les \omterm{def:fm:compliance-and-market-abuse:pad}{transactions personnelles}, puis faire passer un flux d'alertes de surveillance dans une file de tri avec un volume fixe d'heures d'analystes
et trouver le seuil qui capte le plus de vrais cas. \textbf{État final~:} les décisions d'\omterm{def:fm:compliance-and-market-abuse:pad}{autorisation préalable} par motif et le graphique de la file
(la \cref{fig:fm:compliance-and-market-abuse:queue}).
\begin{tutsteps}
\item \textbf{Les alertes.} Les données synthétiques de spoofing du livre~9, \texttt{firm.surveil.spoofing\_days}~: 11\,600 jours-compte, dont 100 abus introduits,
notés par son meilleur détecteur (écart d'exécution des grands ordres multiplié par les annulations après l'exécution de l'autre côté).
\item \textbf{Les analystes.} Chaque alerte prend une demi-heure~; l'équipe dispose de 50 heures pour ce détecteur (100 alertes), ou de 150 (300 alertes).
\item \textbf{La file.} \texttt{fm\_compliance.queue(hours)} balaie le taux d'alerte de 0.1~\% à 20~\% avec \texttt{firm.compliance.triage}.
\item \textbf{Les demandes.} \texttt{fm\_compliance.year()} envoie 1\,215 demandes de \omterm{def:fm:compliance-and-market-abuse:pad}{transactions personnelles} émanant de 150 salariés via
\texttt{firm.compliance.preclear} face à des listes qui changent au fil de l'année.
\end{tutsteps}
\end{tutorial}

\begin{omfigure}
\begin{tikzpicture}
\begin{axis}[width=10.5cm,height=5.4cm,xmode=log,log basis x=10,xtick={0.1,0.2,0.5,1,2,5,10,20},xticklabels={0.1,0.2,0.5,1,2,5,10,20},xlabel={\small taux d'alerte (\% des jours-compte, échelle logarithmique)},
  ylabel={\small vrais cas trouvés (sur 100)},xmin=0.1,xmax=20,ymin=0,ymax=100,
  tick label style={font=\footnotesize},grid=major,grid style={gray!20},table/col sep=comma,
  legend cell align=left,legend style={font=\scriptsize,at={(0.5,-0.3)},anchor=north,/tikz/every even column/.append style={column sep=8pt}},legend columns=2]
\addplot[omDef,thick,mark=*,mark size=1pt] table[x=rate,y=h50]{figdata/desk/16-compliance-and-market-abuse/queue.csv};
\addplot[omThm,thick,dashed,mark=square*,mark size=1pt] table[x=rate,y=h150]{figdata/desk/16-compliance-and-market-abuse/queue.csv};
\draw[gray,dotted] (axis cs:5,0) -- (axis cs:5,100);
  \node[font=\scriptsize,anchor=south west] at (axis cs:5.3,82) {par défaut 5~\%};
\legend{50 heures d'analystes,150 heures d'analystes}
\end{axis}
\end{tikzpicture}

\omcaption{Vrais cas trouvés par une file d'alertes examinée dans l'ordre d'arrivée, en fonction du taux d'alerte, pour deux budgets d'analystes. Chaque courbe culmine là où les
alertes remplissent tout juste la capacité (la \cref{prop:fm:compliance-and-market-abuse:queue})~; au taux par défaut de 5~\% la file déborde et la plupart des vraies alertes
ne sont jamais lues. Données~: \texttt{fm\_compliance.queue}.}\label{fig:fm:compliance-and-market-abuse:queue}
\end{omfigure}

Avec 50 heures d'analystes, le meilleur taux d'alerte est de 0.8~\%, juste en dessous du taux de capacité de $100/11\,600=0.86~\%$, et la file trouve 65 des 100
cas introduits. Au taux par défaut de 5~\%, le détecteur émet des alertes sur 580 jours-compte, dont 89 sont vraies, mais les analystes qui lisent les 100 premières dans
l'ordre d'arrivée en trouvent 15 en moyenne~: 85~\% des abus passent inaperçus, non pas à cause du détecteur, qui en a signalé la plus grande part, mais à cause de la file
(la \cref{fig:fm:compliance-and-market-abuse:funnel}). Tripler les heures déplace le meilleur taux à 2.5~\% et porte à 81 le nombre de vrais cas trouvés~; au taux par défaut
la file en trouve 46. Examiner dans l'ordre des scores plutôt que dans l'ordre d'arrivée rattrape ce que le seuil fait perdre~: 66 trouvés au taux par défaut avec 50 heures,
81 avec 150. Un détecteur naïf, le ratio ordres sur transactions, n'en trouve aucun parmi les 100, quel que soit le seuil, parce que les teneurs de marché du livre~9 annulent davantage
que tout spoofer.

\begin{omfigure}
\begin{tikzpicture}[font=\footnotesize,>=stealth,box/.style={draw,rounded corners,align=center,minimum height=0.8cm}]
\node[box,fill=gray!15,text width=10cm] (a) at (0,3.0) {11\,600 jours-compte évalués, 100 cas introduits};
\node[box,fill=omDef!15,text width=7.6cm] (b) at (0,1.8) {580 alertes au taux par défaut de 5~\%, dont 89 vraies};
\node[box,fill=omProp!15,text width=5cm] (c) at (0,0.6) {100 examinées en 50 heures};
\node[box,fill=omThm!15,text width=2.8cm] (d) at (0,-0.6) {15 vraies trouvées};
\draw[->] (a) -- (b); \draw[->] (b) -- (c); \draw[->] (c) -- (d);
\node[font=\scriptsize,align=left,anchor=west] at (3.9,1.8) {détecteur~: rappel de 89~\%};
\node[font=\scriptsize,align=left,anchor=west] at (2.7,0.6) {file~: 17~\% des alertes lues};
\end{tikzpicture}

\omcaption{L'entonnoir de surveillance au seuil par défaut~: le détecteur trouve la plus grande part des abus, et la file les laisse passer. Données~:
\texttt{fm\_compliance.queue}.}\label{fig:fm:compliance-and-market-abuse:funnel}
\end{omfigure}

\section{Transactions personnelles}

Les employés qui savent ce que la société s'apprête à négocier, ou ce qu'un client va annoncer, peuvent en tirer profit sur leurs propres comptes. Ces règles sont
plus anciennes que le trading algorithmique et s'appliquent à toute personne ayant accès à l'information.

\begin{definition}[Transactions personnelles, autorisation préalable]\label{def:fm:compliance-and-market-abuse:pad}
Les \emph{transactions personnelles}\index{transactions personnelles} (personal account dealing) désignent les opérations réalisées par les employés sur leurs propres comptes ou sur des comptes qu'ils contrôlent ou
dont ils tirent profit. L'\emph{autorisation préalable}\index{autorisation préalable} (pre-clearance) est l'exigence selon laquelle un employé obtient l'approbation de la fonction conformité avant
une opération personnelle. Elle est refusée si l'instrument figure sur une \omterm{def:fm:compliance-and-market-abuse:lists}{liste d'interdiction}, si la société négocie ou s'apprête à négocier cet instrument, si une durée minimale de détention
n'est pas écoulée, ou si l'employé détient une \omterm{def:fm:compliance-and-market-abuse:abuse}{information privilégiée} à son sujet.
\end{definition}

\begin{definition}[Liste d'interdiction, liste de surveillance]\label{def:fm:compliance-and-market-abuse:lists}
Une \emph{liste d'interdiction}\index{liste d'interdiction} (restricted list) recense les instruments dans lesquels la société et ses employés ne peuvent pas négocier, ou ne peuvent négocier qu'avec
une autorisation, généralement parce que la société détient une \omterm{def:fm:compliance-and-market-abuse:abuse}{information privilégiée} sur l'émetteur et le dit publiquement ou à son personnel. Une \emph{liste de surveillance}\index{liste de surveillance} (watch
list) recense, de façon confidentielle et à l'usage de la seule conformité, les instruments au sujet desquels une partie de la société détient une \omterm{def:fm:compliance-and-market-abuse:abuse}{information privilégiée},
afin que les opérations sur ces instruments, où qu'elles aient lieu dans la société, puissent être surveillées sans révéler l'information.
\end{definition}

\omcode{code/firm/compliance/firm_compliance.py}{41}{57}{Autorisation préalable d'une opération personnelle~: la première règle applicable donne le motif
d'un refus.}\label{lst:fm:compliance-and-market-abuse:preclear}

\begin{omfigure}
\begin{tikzpicture}
\begin{axis}[width=10.5cm,height=4.8cm,xbar,bar width=9pt,xmin=0,xmax=1100,xlabel={\small demandes},ytick={0,1,2,3,4,5},
  yticklabels={approuvées,société a négocié récemment,liste d'interdiction,initié,période de détention,liste de surveillance},y dir=reverse,ymin=-0.6,ymax=5.6,
  tick label style={font=\footnotesize},grid=major,grid style={gray!20},table/col sep=comma,nodes near coords,
  nodes near coords style={font=\scriptsize}]
\addplot[fill=omDef!55,draw=omDef] table[y=pos,x=count]{figdata/desk/16-compliance-and-market-abuse/reasons.csv};
\end{axis}
\end{tikzpicture}

\omcaption{Une année de demandes d'opérations personnelles par décision~: 1\,215 demandes, 964 approuvées~; quatre refus sur cinq concernent des noms que la société a négociés
dans les deux jours précédents. Données~: \texttt{fm\_compliance.reasons}.}\label{fig:fm:compliance-and-market-abuse:reasons}
\end{omfigure}

Sur les 1\,215 demandes de l'année, 964 (79.3~\%) ont été approuvées. Sur les 251 refus, 202 (80.5~\%) concernaient des noms
que la société avait négociés dans les deux jours précédents, ce qui relève de la règle anti-front-running~; 23 concernaient des noms de la \omterm{def:fm:compliance-and-market-abuse:lists}{liste d'interdiction}, 5 des noms de la \omterm{def:fm:compliance-and-market-abuse:lists}{liste de surveillance}, 6 des cessions pendant la période de détention de 30 jours, et 15
étaient des demandes d'employés ayant franchi la muraille pour les noms de leurs projets, et chacune des demandes introduites a été refusée. Les refus liés à la \omterm{def:fm:compliance-and-market-abuse:lists}{liste de surveillance} montrent
le coût de cette liste~: un employé refusé pour un motif que la conformité ne peut pas donner apprend qu'il se passe quelque chose sur ce nom. Une société peut réduire ce
signal en ne donnant aucun motif pour aucun refus, ou en surveillant les noms de la \omterm{def:fm:compliance-and-market-abuse:lists}{liste de surveillance} au lieu de refuser les opérations sur ces noms.

\section{Barrières à l'information, listes d'initiés et franchissement de la muraille}

Une société qui conseille des émetteurs, leur prête des fonds ou souscrit leurs titres détient une \omterm{def:fm:compliance-and-market-abuse:abuse}{information privilégiée} dans l'une de ses parties et négocie dans une autre. Elle les tient
séparées au moyen de barrières, et elle consigne qui les franchit.

\begin{definition}[Barrière à l'information, franchissement de la muraille, liste d'initiés]\label{def:fm:compliance-and-market-abuse:barrier}
Une \emph{barrière à l'information}\index{barrière à l'information} (muraille de Chine) est l'ensemble des contrôles organisationnels, physiques et
électroniques qui empêchent l'\omterm{def:fm:compliance-and-market-abuse:abuse}{information privilégiée} détenue dans une partie d'une société d'atteindre les personnes d'une autre partie qui négocient ou conseillent sur les instruments concernés. Un
\emph{franchissement de la muraille}\index{franchissement de la muraille} (wall-crossing) est la transmission contrôlée d'une \omterm{def:fm:compliance-and-market-abuse:abuse}{information privilégiée} à une personne située de l'autre côté d'une barrière, ou à un
investisseur extérieur sollicité au sujet d'une opération, avec son consentement préalable, consigné, et l'engageant à ne pas négocier tant que l'information n'est pas publique ou n'est plus
privilégiée. Une \emph{liste d'initiés}\index{liste d'initiés} (insider list) recense chaque personne ayant accès à une \omterm{def:fm:compliance-and-market-abuse:abuse}{information privilégiée} sur un émetteur, la date à laquelle
elle l'a obtenue et la raison de cet accès.
\end{definition}

\begin{omfigure}
\begin{tikzpicture}[font=\footnotesize,>=stealth,box/.style={draw,rounded corners,align=center,minimum height=0.8cm,text width=3cm}]
\node[box,fill=omDef!15] (priv) at (-3.8,0) {côté privé\\{\scriptsize conseil, prêt, marchés de capitaux}};
\node[box,fill=omProp!15] (pub) at (3.8,0) {côté public\\{\scriptsize négociation, recherche, ventes}};
\draw[very thick,gray] (0,-1.4) -- (0,1.6) node[above,font=\scriptsize,text=black] {barrière à l'information};
\node[box,fill=omThm!15,text width=3.4cm] (comp) at (0,-2.6) {conformité\\{\scriptsize liste de surveillance, liste d'initiés, liste d'interdiction}};
\draw[->,dashed] (priv) -- (pub); \node[font=\scriptsize,align=center] at (-1.15,0.42) {franchissement\\de la muraille};
\draw[->] (priv.south) |- (comp.west); \draw[->] (comp.east) -| (pub.south);
\node[font=\scriptsize,align=left,anchor=west] at (3.95,-1.3) {autorisation préalable,\\surveillance};
\node[font=\scriptsize,anchor=east] at (-3.95,-1.3) {opérations consignées};
\end{tikzpicture}

\omcaption{\omterm{def:fm:compliance-and-market-abuse:barrier}{Barrières à l'information}~: le côté privé détient l'\omterm{def:fm:compliance-and-market-abuse:abuse}{information privilégiée}, le côté public négocie~; la conformité se trouve à cheval sur les deux, tient les listes,
consigne les \omterm{def:fm:compliance-and-market-abuse:barrier}{franchissements de la muraille} et soumet le côté public à l'\omterm{def:fm:compliance-and-market-abuse:pad}{autorisation préalable} et à la surveillance. Schéma.}\label{fig:fm:compliance-and-market-abuse:barrier}
\end{omfigure}

Une personne ayant franchi la muraille figure sur la \omterm{def:fm:compliance-and-market-abuse:barrier}{liste d'initiés} de l'émetteur du projet dès le moment du franchissement, jusqu'à ce que l'information soit rendue publique ou perde
sa pertinence (le projet est alors purgé)~; l'\omterm{def:fm:compliance-and-market-abuse:pad}{autorisation préalable} refuse ses opérations sur ce nom pendant toute cette période, et la surveillance suit ce nom dans
l'ensemble de la société. Les registres remplissent deux fonctions~: ils rendent les contrôles efficaces, et ce sont eux que la société présente à son régulateur lorsqu'une opération sur ce nom est contestée.
Une consultation de marché menée selon les règles du règlement est le même acte accompli auprès d'un investisseur extérieur, avec le même consentement et les mêmes registres.

\begin{method}[Faire fonctionner la fonction conformité]\label{met:fm:compliance-and-market-abuse:run}
\begin{enumerate}
\item Cartographier les activités de la société selon les abus que chacune pourrait commettre ou faciliter, et associer chacune à un détecteur~; vérifier la couverture chaque
fois qu'on ajoute une activité, un produit ou une plateforme.
\item Fixer le seuil de chaque détecteur de façon que ses alertes remplissent les heures des analystes, et examiner dans l'ordre des scores~; mesurer le taux d'alertes vraies et ajuster.
\item Tenir les listes datées~; soumettre chaque opération personnelle à l'\omterm{def:fm:compliance-and-market-abuse:pad}{autorisation préalable}~; consigner chaque \omterm{def:fm:compliance-and-market-abuse:barrier}{franchissement de la muraille} et chaque purge.
\item Déclarer sans délai les ordres et les transactions suspects, et consigner les motifs de chaque décision de ne pas déclarer.
\item Tester la chaîne complète une fois par an avec des cas introduits, comme le fait le tutoriel.
\end{enumerate}
\end{method}

\section{Construction~: le moteur de conformité}\label{bld:fm:compliance-and-market-abuse:compliance}

\begin{build}
\bfield{Rôle} \omterm{def:fm:compliance-and-market-abuse:lists}{Listes d'interdiction} et de surveillance au fil du temps, \omterm{def:fm:compliance-and-market-abuse:pad}{autorisation préalable} des opérations personnelles, un registre des \omterm{def:fm:compliance-and-market-abuse:barrier}{listes d'initiés} avec \omterm{def:fm:compliance-and-market-abuse:barrier}{franchissements de la muraille}, et la
capacité d'une file d'alertes sur les détecteurs de firm.surveil.
\bfield{Interface} \texttt{firm.compliance}~: \texttt{Lists}, \texttt{Rules}, \texttt{preclear}, \texttt{Register} (\texttt{cross},
\texttt{cleanse}, \texttt{insiders}), \texttt{triage}, \texttt{best\_rate}.
\bfield{Règles} La première règle applicable décide d'un refus et est consignée~; les listes et les registres sont datés~; les initiés d'un projet purgé quittent la
liste le jour de la purge~; le tri n'examine jamais plus d'alertes que la capacité.
\bfield{Tests d'acceptation} \texttt{code/firm/compliance/tests/}~: chaque motif de refus~; le registre à travers un projet~; la file sur un détecteur
de précision connue.
\bfield{Pour aller plus loin} Alertes de plusieurs détecteurs partageant une même équipe~; durées de revue qui dépendent de l'alerte~; le délai de déclaration auprès du régulateur comme
contrainte sur la file.
\end{build}

\begin{omsources}
\item Financial Conduct Authority, communiqué de presse du 19 août 2022 concernant Citigroup Global Markets Limited.
\item Regulation (EU) No 596/2014, articles 7, 8, 11, 14--16 et 18.
\item 17 CFR 240.10b-5~; 15 U.S.C. 78o(g)~; 17 CFR 275.204A-1.
\end{omsources}

\section{Exercices}

\begin{exercise}[$\star$]\label{exo:fm:compliance-and-market-abuse:1}
Laquelle des règles de refus du chapitre bloquerait l'achat, par un trader, d'une action que la firme a achetée la veille, et pourquoi cette règle existe-t-elle~?
\end{exercise}

\begin{exercise}[$\star$]\label{exo:fm:compliance-and-market-abuse:2}
Un détecteur attribue un score à 11\,600 jours-compte et les analystes peuvent examiner 100 alertes. Quel taux d'alerte remplit la file~?
\end{exercise}

\begin{exercise}[$\star$]\label{exo:fm:compliance-and-market-abuse:3}
Quelle est la différence entre une \omterm{def:fm:compliance-and-market-abuse:lists}{liste d'interdiction} et une \omterm{def:fm:compliance-and-market-abuse:lists}{liste de surveillance}, et pourquoi la \omterm{def:fm:compliance-and-market-abuse:lists}{liste de surveillance} est-elle confidentielle~?
\end{exercise}

\begin{exercise}[$\star\star$]\label{exo:fm:compliance-and-market-abuse:4}
Au taux par défaut de 5~\% le détecteur signale 89 des 100 cas. Pourquoi la file n'en retrouve-t-elle que 15~?
\end{exercise}

\begin{exercise}[$\star\star$]\label{exo:fm:compliance-and-market-abuse:5}
Pourquoi le ratio ordres sur transactions ne repère-t-il aucun des cas de spoofing placés~?
\end{exercise}

\begin{exercise}[$\star\star$]\label{exo:fm:compliance-and-market-abuse:6}
Un analyste de recherche franchit la muraille pour une OPA le jour 40, et le projet est purgé le jour 75. Quels jours une transaction personnelle sur la cible
est-elle refusée, et pourquoi la surveillance de l'ensemble des opérations de la firme sur ce titre reste-t-elle nécessaire~?
\end{exercise}

\begin{exercise}[$\star\star\star$]\label{exo:fm:compliance-and-market-abuse:7}
\emph{Programmation.} Exécuter \texttt{fm\_compliance.queue} avec 150 heures, dans l'ordre d'arrivée et dans l'ordre des scores. Que trouvent les meilleurs taux et le taux
par défaut~?
\end{exercise}

\begin{exercise}[$\star\star\star$]\label{exo:fm:compliance-and-market-abuse:8}
\emph{Trouver la faille.} «~Nous avons abaissé chaque seuil pour attraper davantage d'abus~; notre surveillance est plus forte que l'année dernière.~»
\end{exercise}

\section{Problème~: la file d'alertes}

\begin{problem}[{Problème du week-end --- la file d'alertes}]\label{pb:fm:compliance-and-market-abuse:1}
Un nouveau responsable de la conformité hérite d'une équipe de surveillance qui vide sa file chaque trimestre en fermant les anciennes alertes sans les lire,
ainsi qu'un processus d'\omterm{def:fm:compliance-and-market-abuse:pad}{autorisation préalable} mené par courriel.

\medskip\noindent\textbf{Partie I --- Les règles.}
\begin{enumerate}
\item Définir l'\omterm{def:fm:compliance-and-market-abuse:abuse}{abus de marché}, l'\omterm{def:fm:compliance-and-market-abuse:abuse}{information privilégiée} et l'\omterm{def:fm:compliance-and-market-abuse:abuse}{opération d'initié}.
\item Qu'exige le règlement européen des firmes qui organisent ou exécutent des transactions~?
\item Énoncer les règles américaines sur le délit d'initié, les politiques des broker-dealers et les codes de déontologie des conseillers.
\item Quelle défaillance la FCA a-t-elle relevée chez Citigroup Global Markets en 2022~?
\end{enumerate}

\medskip\noindent\textbf{Partie II --- La file.}
\begin{enumerate}[resume]
\item Définir une \omterm{def:fm:compliance-and-market-abuse:stor}{déclaration d'opérations et d'ordres suspects} (STOR).
\item Énoncer et démontrer la \cref{prop:fm:compliance-and-market-abuse:queue}.
\item Donner les meilleurs taux et les cas véritables trouvés pour 50 et 150 heures d'analyste.
\item Donner ce que le taux par défaut trouve et manque avec 50 heures.
\item Qu'est-ce que la revue dans l'ordre des scores change~?
\item Pourquoi le rappel d'un détecteur n'est-il pas le rappel de la fonction~?
\end{enumerate}

\medskip\noindent\textbf{Partie III --- \omterm{def:fm:compliance-and-market-abuse:pad}{Transactions personnelles} et barrières.}
\begin{enumerate}[resume]
\item Définir les \omterm{def:fm:compliance-and-market-abuse:pad}{transactions personnelles}, l'\omterm{def:fm:compliance-and-market-abuse:pad}{autorisation préalable}, la \omterm{def:fm:compliance-and-market-abuse:lists}{liste d'interdiction} et la \omterm{def:fm:compliance-and-market-abuse:lists}{liste de surveillance}.
\item Donner les décisions d'\omterm{def:fm:compliance-and-market-abuse:pad}{autorisation préalable} de l'année, par motif.
\item Définir une \omterm{def:fm:compliance-and-market-abuse:barrier}{barrière à l'information}, un \omterm{def:fm:compliance-and-market-abuse:barrier}{franchissement de la muraille} et une \omterm{def:fm:compliance-and-market-abuse:barrier}{liste d'initiés}.
\item Combien de temps une personne ayant franchi la muraille reste-t-elle sur la \omterm{def:fm:compliance-and-market-abuse:barrier}{liste d'initiés}, et qu'en est-il de ses opérations~?
\end{enumerate}

\medskip\noindent\textbf{Partie IV --- Le plan.}
\begin{enumerate}[resume]
\item Qu'y a-t-il de problématique dans la fermeture des anciennes alertes sans les lire~?
\item Comment fixer le seuil de chaque détecteur~?
\item Comment vérifier la couverture lorsque la firme commence à négocier un nouveau produit~?
\item Comment tester l'ensemble de la chaîne~?
\item Énoncer le \emph{résultat nommé}~: le seuil d'alerte qui maximise le nombre de cas véritables trouvés pour un nombre fixe d'heures d'analyste, et la part
manquée au seuil par défaut.
\item En deux phrases, rédiger le plan.
\end{enumerate}
\end{problem}

\section{Questions d'entretien}

\begin{interviewq}[$\star$ \iqroles{trader}]\label{iq:fm:compliance-and-market-abuse:1}
Qu'est-ce qu'une \omterm{def:fm:compliance-and-market-abuse:abuse}{information privilégiée}, et que devez-vous faire si vous en recevez une par erreur~?
\end{interviewq}

\begin{interviewq}[$\star$ \iqroles{trader, chercheur}]\label{iq:fm:compliance-and-market-abuse:2}
Pourquoi votre firme vous oblige-t-elle à soumettre vos \omterm{def:fm:compliance-and-market-abuse:pad}{transactions personnelles} à une \omterm{def:fm:compliance-and-market-abuse:pad}{autorisation préalable}~?
\end{interviewq}

\begin{interviewq}[$\star\star$ \iqroles{développeur}]\label{iq:fm:compliance-and-market-abuse:3}
Concevez le modèle de données d'une \omterm{def:fm:compliance-and-market-abuse:lists}{liste d'interdiction} capable de répondre à la question «~ce titre était-il sous restriction le 3 mars à 10~h~32~?~».
\end{interviewq}

\begin{interviewq}[$\star\star$ \iqroles{chercheur}]\label{iq:fm:compliance-and-market-abuse:4}
Un détecteur de surveillance a un rappel de 90~\% pour un taux d'alerte de 5~\%. Est-il bon~?
\end{interviewq}

\begin{interviewq}[$\star\star$ \iqroles{risque}]\label{iq:fm:compliance-and-market-abuse:5}
Qu'est-ce qu'un \omterm{def:fm:compliance-and-market-abuse:barrier}{franchissement de la muraille}, et quelles traces doit-il laisser~?
\end{interviewq}

\begin{interviewq}[$\star\star\star$ \iqroles{chercheur, développeur}]\label{iq:fm:compliance-and-market-abuse:6}
Plusieurs détecteurs partagent une même équipe d'analystes. Comment répartiriez-vous leurs heures~?
\end{interviewq}
